% %tabulacao para a listagem
% %\newcommand{\itab}[1]{\hspace{0em}\rlap{#1}}
% %\newcommand{\tab}[1]{\hspace{.2\textwidth}\rlap{#1}}
% %inicio do capitulo



% %Cap 1 - Introdução 
% %   --> Tentem escrever nessa ordem (não precisa criar subseção): 
% %        Contextualização do Tema; Descrição do Problema; Proposta; Limitações; Justificativa
% %   1.1 Objetivos: Geral e Específicos
% %   1.2 Materiais e Métodos

% \chapter[INTRODUÇÃO]{INTRODUÇÃO}
% O câncer de pele é amplamente prevalente em nível global, representando uma grande parcela dos diagnósticos oncológicos em países desenvolvidos e em desenvolvimento. Segundo estimativas de \cite{cp2023inca}, mais de 30\% dos casos de câncer no Brasil no triênio de 2023-2025 serão de câncer de pele. De acordo com a Sociedade Brasileira de Dermatologia \cite{SBD_2024}, os tipos mais comuns são o carcinoma basocelular e o carcinoma espinocelular, enquanto o melanoma, embora mais raro, é a variante mais agressiva e com maior taxa de mortalidade. A incidência de câncer de pele tem aumentado em países com alta exposição solar, como o Brasil, devido à intensa radiação ultravioleta (UV). Estudos indicam que a falta de proteção solar é um fator determinante para esse aumento.

% Além da exposição à radiação UV, fatores como predisposição genética, pigmentação cutânea, histórico de queimaduras solares e idade têm impacto relevante no desenvolvimento do câncer de pele \cite{shraim2023systematic}. No contexto brasileiro, a conscientização sobre o diagnóstico precoce é fundamental, especialmente em regiões com acesso limitado a especialistas. A detecção precoce é crucial para reduzir a mortalidade e os custos com o tratamento de estágios avançados. Nesse cenário, tecnologias baseadas em inteligência artificial (IA) surgem como uma alternativa promissora, ampliando o acesso ao diagnóstico e aumentando a precisão, mesmo em locais com escassez de especialistas.

% Esta pesquisa propõe o desenvolvimento de um modelo baseado na rede neural LAMMA3, treinado, validado e testado com imagens fornecidas pela Society for Imaging Informatics in Medicine (SIIM) e pelo International Skin Imaging Collaboration (ISIC). O objetivo é maximizar a acurácia no diagnóstico e reduzir a incidência de falsos positivos em lesões cutâneas suspeitas. A proposta acompanha os recentes avanços do aprendizado profundo na medicina, buscando oferecer uma solução inovadora que una precisão diagnóstica e acessibilidade, atendendo às crescentes demandas por tecnologias eficazes e inclusivas na dermatologia.

% \section{Motivação}
% Nos últimos anos, as técnicas de aprendizado profundo e aprendizado de máquina têm demonstrado grande potencial como ferramentas auxiliares em diagnósticos médicos, especialmente em áreas com escassez de especialistas. Na dermatologia, esses métodos têm se mostrado essenciais para a detecção precoce de lesões cutâneas, permitindo que sistemas baseados em IA reconheçam padrões em imagens de pele e identifiquem sinais de malignidade com acurácia comparável à de especialistas. Devido à sua extensão territorial e desigualdades regionais, o Brasil enfrenta desafios significativos para garantir o diagnóstico precoce para toda a população. Assim, o desenvolvimento de modelos computacionais aplicáveis remotamente surge como uma alternativa viável para promover uma saúde pública mais acessível.

% Além disso, a eficácia de algoritmos no diagnóstico de melanoma é cada vez mais relevante, visto que o melanoma, embora menos comum, é responsável pela maior taxa de mortalidade entre os cânceres de pele. Esta pesquisa visa demonstrar a capacidade de um modelo avançado de redes neurais para melhorar a detecção precoce do câncer de pele e, assim, contribuir para aumentar a sobrevida de pacientes diagnosticados com a doença. A presente investigação se destaca pela integração de redes neurais para a segmentação de lesões cutâneas, complementadas por sistemas multiagentes para análise e diagnóstico, permitindo maior confiabilidade na tomada de decisões.

% \section{Lacuna do Conhecimento}
% Embora existam modelos que realizam com precisão a segmentação de imagens de lesões cutâneas e a classificação das mesmas, ainda há escassez de pesquisas que explorem a combinação desses resultados com sistemas multiagentes para diagnósticos automatizados e precisos. A integração dessas abordagens é uma área pouco explorada, e a combinação de segmentação, classificação e sistemas multiagentes pode melhorar o diagnóstico diferencial em câncer de pele.

% A literatura atual foca amplamente na segmentação e classificação de lesões de pele, mas observa-se uma limitação na aplicação prática desses algoritmos em contextos clínicos, devido à falta de integração com sistemas que gerem diagnósticos em tempo real a partir de dados coletados automaticamente. Ao implementar um sistema multiagente em conjunto com redes neurais, espera-se que esta pesquisa ofereça uma contribuição inédita, preenchendo a lacuna de integração entre processamento de imagens e suporte ao diagnóstico automatizado. Essa inovação visa proporcionar um diagnóstico mais preciso e viável em regiões com acesso limitado a especialistas, contribuindo para a democratização do diagnóstico de câncer de pele.

% \section{Proposta da Pesquisa}
% Esta pesquisa propõe a aplicação de sistemas multiagentes para desenvolver um diagnóstico diferencial, avaliando a acurácia e a eficiência diagnóstica do sistema com base em imagens segmentadas e classificações externas.

% A arquitetura proposta será organizada em três módulos principais: (i) segmentação de lesões cutâneas, (ii) classificação da natureza da lesão (benigna ou maligna), e (iii) diagnóstico diferencial via sistema multiagente com validação crítica. O desempenho será avaliado por meio de métricas clínicas como acurácia, sensibilidade e precisão, bem como por critérios de completude e consistência textual dos pareceres gerados.

% Apesar do potencial da abordagem proposta, este trabalho apresenta limitações inerentes ao uso de modelos de linguagem para diagnóstico clínico, incluindo a sensibilidade a variações de prompt, a ausência de contexto clínico complementar e a dependência da qualidade das segmentações. Além disso, a validação do sistema ocorre em um ambiente experimental controlado, o que pode limitar a generalização para contextos clínicos reais. Tais limitações são reconhecidas e apontam para direções futuras de pesquisa.

% Esta pesquisa propõe a aplicação de sistemas multiagentes para desenvolver um diagnóstico diferencial, avaliando a acurácia e a eficiência diagnóstica do sistema com base em imagens segmentadas e classificações externas.

% \section{Objetivos}
% \subsection{Objetivo Principal}
% Avaliar a eficácia de modelos multiagentes no diagnóstico diferencial do câncer de pele, em combinação com algoritmos de segmentação e classificação de imagem.

% \subsection{Objetivos Específicos}
% \begin{itemize}
%     \item Realizar revisão bibliográfica sobre o tema;
%     \item Desenvolver um algoritmo de segmentação para imagens de lesões de pele;
%     \item Desenvolver um algoritmo de classificação para categorizar as lesões conforme o tipo (ou ausência) de câncer de pele;
%     \item Implementar um sistema multiagente capaz de realizar o diagnóstico diferencial.
% \end{itemize}


\chapter[INTRODUÇÃO]{INTRODUÇÃO}

O câncer de pele é uma das neoplasias mais prevalentes mundialmente, representando parcela significativa dos diagnósticos oncológicos em países desenvolvidos e em desenvolvimento. Estimativas do Instituto Nacional de Câncer (INCA) indicam que, no Brasil, mais de 30\% dos casos de câncer previstos para o triênio 2023–2025 serão de câncer de pele \cite{cp2023inca}. De acordo com a Sociedade Brasileira de Dermatologia (SBD) \cite{SBD_2024}, os tipos mais comuns são o carcinoma basocelular e o carcinoma espinocelular, enquanto o melanoma, embora menos frequente, apresenta alta letalidade e maior taxa de mortalidade. Em países com elevada exposição solar, como o Brasil, a incidência tende a aumentar devido à intensa radiação ultravioleta (UV), sendo a ausência de proteção solar um fator determinante nesse cenário.

Além da radiação UV, fatores como predisposição genética, pigmentação cutânea, histórico de queimaduras solares e idade influenciam significativamente o desenvolvimento da doença \cite{shraim2023systematic}. No contexto brasileiro, a conscientização sobre diagnóstico precoce é essencial, sobretudo em áreas com acesso limitado a especialistas. A detecção precoce reduz mortalidade e custos associados ao tratamento em estágios avançados. Nesse sentido, tecnologias baseadas em inteligência artificial (IA) mostram-se promissoras para ampliar o acesso ao diagnóstico e elevar a precisão, mesmo em regiões com escassez de recursos humanos especializados.

Esta pesquisa propõe o desenvolvimento de um sistema baseado em \textit{Large Language Models} (LLMs), mais especificamente os modelos GPT-5.4-nano \cite{openai2025gpt5nano} (OpenAI) e Gemini-3.1-flash-lite \cite{google2025gemini31} (Google), empregando sua capacidade multimodal de análise visual e geração de linguagem natural para a \textbf{detecção de melanoma} em imagens dermoscópicas. Os experimentos utilizam o conjunto PH2 \cite{mendonca2013ph2}, composto por 200 imagens dermoscópicas com diagnósticos e máscaras de segmentação anotadas por especialistas, o que permite isolar a avaliação diagnóstica das LLMs da variabilidade introduzida por modelos de segmentação automática. O objetivo central é avaliar a acurácia diagnóstica, a interpretabilidade clínica e a consistência textual dos pareceres gerados, contribuindo para o desenvolvimento de ferramentas de triagem assistida por IA acessíveis e explicáveis.

\section{Motivação}

Nos últimos anos, técnicas de aprendizado profundo e aprendizado de máquina vêm demonstrando alto potencial como ferramentas auxiliares em diagnósticos médicos, com especial utilidade em regiões com déficit de especialistas. Na dermatologia, Esteva et al. \cite{esteva2017dermatologist} demonstraram que redes neurais profundas atingem acurácia equivalente à de dermatologistas na classificação de câncer de pele; Haenssle et al. \cite{haenssle2018man} compararam uma rede convolucional a 58 especialistas em melanoma, obtendo desempenho diagnóstico superior; e Brinker et al. \cite{brinker2019deep} superaram 136 dos 157 dermatologistas avaliados em tarefa de classificação dermatoscópica.
 

O Brasil, devido à sua extensão territorial e desigualdades regionais, enfrenta desafios para garantir diagnóstico precoce universal. Assim, modelos computacionais aplicáveis de forma remota tornam-se alternativas viáveis para promover saúde pública mais acessível. A relevância é ainda maior para o melanoma, que, apesar de menos incidente, responde pela maioria das mortes por câncer de pele. Esta investigação se diferencia pela integração entre modelos de linguagem multimodal com máscaras de segmentação de referência e sistemas multiagentes para análise e emissão de diagnósticos, ampliando a interpretabilidade e a confiabilidade na tomada de decisão.

\section{Lacuna do Conhecimento}

Embora existam modelos capazes de segmentar e classificar imagens de lesões cutâneas com alta precisão, há uma escassez de trabalhos que integrem esses resultados a sistemas multiagentes para diagnósticos automatizados. A combinação de segmentação, LLMs multimodais e suporte multiagente é uma abordagem pouco explorada e com potencial para melhorar a detecção de melanoma de forma automatizada e explicável. Estudos recentes têm demonstrado a eficácia de arquiteturas que unificam segmentação e classificação, como o MT-TransUNet, que emprega \textit{transformers} multitarefa para integrar ambas as etapas \cite{chen2021mt-transunet}, o modelo de \textit{bootstrapping} mútuo (MB-DCNN), que permite que segmentação e classificação se reforcem reciprocamente \cite{xia_mutual_bootstrapping_2019}, e o GS-TransUNet, que combina \textit{Gaussian Splatting} e \textit{Transformer UNet} para alcançar desempenho superior em benchmarks dermatológicos \cite{anand2025gs-transunet}. Ainda assim, poucos estudos avançam além da integração em redes unificadas, propondo arquiteturas baseadas em sistemas multiagentes que possam distribuir e orquestrar diferentes competências diagnósticas de forma colaborativa.


Grande parte da literatura concentra-se nas etapas isoladas de segmentação e classificação, mas carece de aplicações práticas que integrem essas fases a sistemas de suporte à decisão em tempo real. Ao propor a integração entre LLMs multimodais e sistemas multiagentes para a detecção de melanoma, esta pesquisa visa preencher essa lacuna, oferecendo um método capaz de proporcionar diagnósticos explicáveis e aplicáveis a regiões com recursos limitados, contribuindo para a democratização do acesso ao diagnóstico.

\section{Proposta da Pesquisa}

Propõe-se a aplicação de sistemas multiagentes para a detecção de melanoma em imagens dermoscópicas, avaliando a acurácia diagnóstica e a consistência textual dos pareceres gerados a partir de imagens e suas respectivas máscaras de segmentação de referência. Além da comparação de acurácia com um \textit{baseline} supervisionado, este trabalho propõe um \textbf{framework de cinco eixos} para a avaliação automatizada da consistência estrutural e lógica dos relatórios diagnósticos gerados pelas LLMs, complementado por uma \textbf{análise de calibração de incerteza} que investiga se os modelos concentram sua hesitação diagnóstica nos casos clinicamente mais ambíguos. Essa dupla contribuição metodológica constitui o principal diferencial desta pesquisa frente à literatura de classificação binária tradicional.

A arquitetura será composta por três módulos principais:
\begin{enumerate}
    \item Uso das máscaras de segmentação de referência do conjunto PH2 como entrada padronizada ao sistema;
    \item Detecção de melanoma por modelos de linguagem multimodal (melanoma vs.\ não-melanoma);
    \item Diagnóstico de melanoma via sistema multiagente, com consolidação dos pareceres pelo agente sumarizador.
\end{enumerate}

A avaliação contemplará métricas clínicas como acurácia, sensibilidade e precisão, bem como critérios de completude e consistência textual dos pareceres. Entre as limitações previstas, destacam-se a sensibilidade dos modelos de linguagem a variações de \textit{prompt}, a ausência de contexto clínico complementar e o tamanho reduzido do conjunto PH2 (200 imagens), que poderá limitar a generalização dos resultados. A validação ocorrerá em ambiente experimental controlado, o que deverá ser considerado na interpretação das métricas obtidas.

\section{Objetivos}

\subsection{Objetivo Principal}
Avaliar a eficácia de modelos multiagentes na detecção de melanoma em imagens dermoscópicas, utilizando máscaras de segmentação de referência do conjunto PH2 e os protocolos clínicos ABCD, Menzies e SPCL como base do raciocínio diagnóstico estruturado.


\subsection{Objetivos Específicos}
\label{sec:objetivos_especificos}
\begin{itemize}
    \item Objetivos até a qualificação
    \begin{itemize}
        \item Investigar o estado da arte em segmentação de imagens, classificação de lesões cutâneas e sistemas diagnósticos baseados em inteligência artificial aplicados à dermatologia
        \begin{itemize}
            \item Capítulo de revisão bibliográfica na dissertação;
            \item Documento de mapeamento de protocolos clínicos (ABCD, Menzies, SPCL);
            \item Arquivo de referências bibliográficas organizado (\texttt{.bib});
            \item Quadro comparativo entre abordagens existentes.
        \end{itemize}
    
        \item Preparar o conjunto PH2 e suas máscaras de segmentação de referência como base para a avaliação diagnóstica
        \begin{itemize}
            \item Pré-processamento e normalização das imagens PH2;
            \item Integração das máscaras de referência (\textit{ground truth}) ao \textit{pipeline} de inferência das LLMs;
            \item Organização do conjunto em subconjuntos para avaliação binária (melanoma vs.\ não-melanoma).
        \end{itemize}
    
        \item Avaliar a eficácia de diferentes modelos de linguagem natural (LLMs) na detecção de melanoma com base na imagem dermoscópica e sua máscara de segmentação de referência
        \begin{itemize}
            \item Relatório comparativo das LLMs testadas (GPT-5.4-nano e Gemini-3.1-flash-lite);
            \item Base de prompts (Chain of Thought, Risk Prompt, etc.);
            \item Tabela de desempenho quantitativo com métricas clássicas;
            \item Discussão qualitativa das respostas geradas pelas LLMs.
        \end{itemize}
    
        \item Implementar um sistema multiagente baseado em LLMs para geração de diagnósticos estruturados, considerando diferentes protocolos clínicos
        \begin{itemize}
            \item Arquitetura funcional implementada com LangChain/LangGraph;
            \item Diagrama do fluxo de agentes e estados do sistema;
            \item Diagnósticos gerados pelos agentes com base em diferentes protocolos;
            \item Código-fonte modularizado e reusável.
        \end{itemize}

        \item Validar o desempenho do sistema proposto por meio de comparação com modelos tradicionais de visão computacional e análise de consistência textual
        \begin{itemize}
            \item Treinamento do \textit{baseline} ResNet-50 diretamente sobre o conjunto PH2 (divisão estratificada 70/30, \textit{fine-tuning} completo com ponderação de classes), eliminando a dependência de transferência de conhecimento entre bases distintas;
            \item Tabelas comparativas com o modelo \textit{baseline} retreinado;
            \item Análise de consistência estrutural e lógica das respostas por meio de cinco eixos de verificação objetiva;
            \item Avaliação da completude e estrutura das respostas em linguagem natural;
            \item Discussão sobre limitações, falsos positivos e interpretabilidade.
        \end{itemize}
    \end{itemize}
    \item Objetivos até a defesa
    \begin{itemize}
        \item Investigar uma variante do sistema multiagente na qual \textit{features} associadas aos critérios dos protocolos clínicos (ABCD, Menzies, SPCL) sejam extraídas por métodos tradicionais de visão computacional e fornecidas como entrada estruturada às LLMs, em complemento ao envio direto da imagem
        \begin{itemize}
            \item Extração de \textit{features} morfológicas (assimetria, bordas, cor, estruturas dermoscópicas) por técnicas clássicas de visão computacional;
            \item Comparação de desempenho entre a abordagem por \textit{features} extraídas e a abordagem atual de envio direto da imagem à LLM;
            \item Discussão sobre o impacto da informação estruturada previamente extraída na qualidade do raciocínio diagnóstico gerado.
        \end{itemize}

        \item Implementar de forma completa o Agente Crítico baseado em \textit{Retrieval-Augmented Generation} (RAG), ainda em estágio preliminar na fase de qualificação, e mensurar isoladamente sua contribuição ao diagnóstico final
        \begin{itemize}
            \item Definição da base de conhecimento biomédica a ser consultada (e.g., PubMed);
            \item Definição do método de recuperação e do modelo de \textit{embeddings} utilizado;
            \item Integração do Agente Crítico à versão do sistema efetivamente avaliada;
            \item Análise comparativa do desempenho e da consistência textual do sistema com e sem o Agente Crítico.
        \end{itemize}

        \item Organizar os resultados obtidos em um protótipo funcional e documentar as análises por meio de relatórios técnicos e artigos científicos
        \begin{itemize}
            \item Protótipo final em ambiente executável (e.g., Colab ou App);
            \item Relatório técnico consolidado com testes e resultados;
            \item Artigo científico no formato IEEE ou Elsevier;
            \item Materiais de apoio à defesa (slides e roteiro).
        \end{itemize}
    \end{itemize}
\end{itemize}

A distribuição temporal dessas atividades ao longo do mestrado é detalhada no cronograma de execução apresentado ao final desta dissertação (Figura~\ref{fig:crono}).

\section{Metodologia}

A pesquisa é de natureza aplicada e experimental, estruturada em quatro etapas interdependentes: (i) revisão do estado da arte em detecção de melanoma e diagnóstico assistido por IA em dermatologia; (ii) preparação do conjunto PH2 \cite{mendonca2013ph2} e integração de suas máscaras de segmentação de referência ao \textit{pipeline} diagnóstico, eliminando a variabilidade de modelos de segmentação automática; (iii) avaliação comparativa dos modelos de linguagem multimodal GPT-5.4-nano \cite{openai2025gpt5nano} e Gemini-3.1-flash-lite \cite{google2025gemini31} para detecção de melanoma e geração de pareceres clínicos estruturados; e (iv) implementação e validação de um sistema multiagente coordenado pelos \textit{frameworks} LangChain e LangGraph. O detalhamento de conjuntos de dados, arquiteturas, estratégias de \textit{prompting} e métricas de avaliação é apresentado no Capítulo~4.



\section{Estrutura da Dissertação}
Esta dissertação está organizada em seis capítulos. O Capítulo~2 apresenta o referencial teórico, discutindo os principais conceitos e abordagens relacionados à segmentação de imagens dermoscópicas, ao uso de modelos de linguagem multimodal na medicina e à aplicação de sistemas multiagentes em diagnósticos automatizados. O Capítulo~3 revisa o estado da arte, posicionando este trabalho em relação às pesquisas existentes em classificação de lesões cutâneas, diagnóstico médico assistido por LLMs e sistemas multiagentes para suporte clínico. O Capítulo~4 descreve os materiais e métodos empregados, incluindo a preparação do conjunto PH2, o uso das máscaras de segmentação de referência, a arquitetura multiagente baseada nos protocolos ABCD, Menzies e SPCL, e as estratégias de avaliação adotadas. O Capítulo~5 apresenta e discute os resultados obtidos, tanto quantitativos quanto qualitativos, com ênfase na comparação entre o \textit{baseline} supervisionado e os modelos de linguagem e na análise da consistência textual dos pareceres gerados. Por fim, o Capítulo~6 reúne as conclusões, destacando as contribuições alcançadas, as limitações encontradas e as perspectivas para estudos futuros.
